% Figure from MATH 590 notes, section 9. Compile with the notes preamble.
\begin{tikzpicture}
\begin{axis}[nfig, width=0.62\textwidth, height=0.42\textwidth, clip=false,
  axis lines=middle, xmin=-1.65, xmax=1.75, ymin=-0.12, ymax=1.55,
  xtick={-1,-0.5,0.5,1}, xticklabels={$-1$,$-\tfrac12$,$\tfrac12$,$1$},
  ytick=\empty, xlabel={}, ylabel={}]
  \fill[black!6] (axis cs:-1.6,0.5) rectangle (axis cs:1.6,1);
  \draw[black!40, dotted] (axis cs:-1,0) -- (axis cs:-1,1) (axis cs:-0.5,0) -- (axis cs:-0.5,0.5)
                          (axis cs:0.5,0) -- (axis cs:0.5,0.5) (axis cs:1,0) -- (axis cs:1,1);
  \addplot[figR, very thick, domain=-1.4:0] {-x};
  \addplot[figB, very thick, domain=0:1.4] {x};
  % C on the y-axis
  \draw[black, line width=2pt] (axis cs:0,0.5) -- (axis cs:0,1);
  \node[font=\scriptsize, anchor=east] at (axis cs:-0.03,0.5) {$\tfrac12$};
  \node[font=\scriptsize, anchor=east] at (axis cs:-0.03,1) {$1$};
  \node[font=\scriptsize, anchor=west] at (axis cs:0.05,0.8) {$C$};
  % preimages on the x-axis
  \draw[figR, line width=2.2pt] (axis cs:-1,0) -- (axis cs:-0.5,0);
  \draw[figB, line width=2.2pt] (axis cs:0.5,0) -- (axis cs:1,0);
  \addplot[only marks, mark=*, mark size=1.5pt, figR] coordinates {(-1,0) (-0.5,0)};
  \addplot[only marks, mark=*, mark size=1.5pt, figB] coordinates {(0.5,0) (1,0)};
  \addplot[only marks, mark=*, mark size=1.8pt, black] coordinates {(0,0)};
  \node[font=\scriptsize, figR, anchor=south east] at (axis cs:-1.36,1.4) {$f_2(x) = -x$};
  \node[font=\scriptsize, figB, anchor=south west] at (axis cs:1.36,1.4) {$f_1(x) = x$};
  \node[font=\scriptsize, figR] at (axis cs:-0.75,0.2) {$f_2^{-1}(C)$};
  \node[font=\scriptsize, figB] at (axis cs:0.75,0.2) {$f_1^{-1}(C)$};
\end{axis}
\end{tikzpicture}
